\chapter{Functional Core, Imperative Shell}
\label{ch:fcis}

Functional Core, Imperative Shell (FCIS) is an architectural discipline for separating \textit{what to decide} from \textit{how to reach the outside world}
\autocite{bernhardt2012boundaries,seemannFcis}.

While I am not sure that this pattern is widely used, and I don't have much experience with it, I personally like it for its potential for simplicity and its ability to isolate the pure logic from the outside world, which makes it \textit{very} easy to test. Your logic code can also be broken down into small functions, generally making your code easier to understand and maintain.

The \textbf{functional core} holds business rules, calculations, and data transformations as pure, deterministic code. The \textbf{imperative shell} performs side effects: receiving input from the outside world (presentation layer, web api, cli, gui, etc.),loading data, calling external systems, saving results, and publishing messages. All these activities are interacting with the outside world.

\medskip

Gary Bernhardt popularised the idea in his \textit{Boundaries} talk; since then it has become a practical way to keep complex logic testable in mainstream object-oriented languages.\autocite{googleFcis2025,ssenseFcis,fritzscheFcis}

FCIS is not necessarily a replacement for layered or hexagonal architecture. It can be incorporated into these architectures, or it can be used standalone. It is a cross-cutting refinement you can apply \textit{inside} a controller, a service class, a package, or a module. Any calculation-heavy code path---whether it lives in a REST endpoint, a JavaFX handler, or an application service---benefits from the same split between pure rules and imperative wiring.


I will, however, try to minimize references to hexagonal architecture, and other architectural styles, to keep the focus on the FCIS pattern itself. You may then extract ideas and apply them to other architectures, if you like.

\medskip


In this chapter we use the \texttt{Abyssal Watch} case study with two running examples:
\begin{itemize}
    \item \textbf{SCN-03 --- Triangulate Kaiju Position:} fuse multiple sensor readings into an estimated area (geometry and data fusion),
    \item \textbf{SCN-04 --- Reconcile Mission Telemetry:} merge partially uplinked observations when a Scout Vessel docks (sync and deduplication).
\end{itemize}

The first example revisits triangulation from Chapter~\ref{ch:hexagonal} through an FCIS lens: we extract pure fusion math and apply it \textbf{without} a dedicated service class at first. The second example introduces a \textbf{new} retrieval-time problem---intermittent uplink during mission closure---that has not appeared in prior chapters.

Both examples highlight where functional programming ideas shine: multi-step transformations where you want fast, deterministic unit tests without mocks or databases.

\medskip

We teach FCIS in two stages. \textbf{Stage one} is the minimal shape: one impure layer at the delivery boundary (a Web API controller or JavaFX button handler) and one pure functional core. \textbf{Stage two}, towards the end of the chapter, adds an explicit application service when the same use case must be triggered from several entry points, and separates delivery from infrastructure.

We begin with functional programming concepts, the \textit{imperative sandwich}, the two-layer starting point, and both examples in the simple form.

\section{Functional Programming Concepts}
\label{sec:fcis_fp_concepts}

Before we apply FCIS, we need a shared vocabulary. The terms below appear throughout this chapter and in much of the functional-programming literature.\autocite{seemannFcis} Given that FCIS is a pattern that is based on functional programming concepts, it is important to understand the concepts first, as these are building blocks of the pattern.

\subsection{Functional Programming}
\label{sec:fcis_functional_programming}

\textbf{Functional programming} treats computation primarily as the evaluation of functions over values. Programs are built by composing functions that take inputs and return outputs, rather than by issuing sequences of commands that mutate shared state. In practice, functional style emphasises:
\begin{itemize}
    \item functions as first-class values (they can be passed, returned, and stored),
    \item expressions over statements where possible,
    \item immutable data structures,
    \item and avoiding shared mutable state across concurrent work.
\end{itemize}
Languages such as Haskell and Clojure are fully functional; Java, C\#, and Python can adopt functional \textit{discipline} in selected areas without becoming functional languages wholesale. FCIS is exactly that: apply functional discipline to the parts of your system that are really about calculation and decision-making. That would be the brain of your system, i.e. the business logic.

\subsection{Imperative Programming}
\label{sec:fcis_imperative_programming}

\textbf{Imperative programming} expresses computation as a sequence of commands that change program state. ``Do this, then that, then update this variable'' is the natural shape. Most business applications are imperative at the edges: read a row from the database, send an HTTP response, write a log line. That is not a flaw---the outside world is stateful and effectful. FCIS does not try to eliminate imperative code; it \textit{contains} it in a thin shell so the core stays easy to reason about.

\subsection{Honest Functions}
\label{sec:fcis_honest_functions}

An \textbf{honest function} is one whose signature tells the truth about what it needs and what it does.\autocite{seemannFcis} If a method is named \texttt{calculateScore(observations)} but secretly reads from a database, consults a global configuration object, or logs to the console, it is \textit{dishonest}: callers cannot understand or test it from the parameters alone. Dishonest functions hide dependencies, which forces integration tests, mocks, and surprises in production.

Honest functions take everything they need as explicit arguments and return explicit results. That honesty is what makes the functional core testable: the test passes data in and asserts on data out, with no hidden collaborators.

\subsection{Pure and Impure Functions}
\label{sec:fcis_pure_impure}

A \textbf{pure function} is deterministic and free of side effects: given the same inputs, it always produces the same outputs, and it does not observe or change anything outside its local scope. An \textbf{impure function} does anything else---reads the clock, queries a database, mutates a field on a shared object, or writes a file.

Pure functions are the building blocks of the functional core. Impure functions belong in the imperative shell. The boundary is pragmatic: a method that only transforms in-memory data structures is pure; a method that opens a JDBC connection is not.

Code snippet~\ref{lst:fcis-impure-summarize} shows an impure method that mixes a database read with a simple calculation. Code snippet~\ref{lst:fcis-pure-summarize} extracts the calculation into a pure function; the shell (not shown here) would load observations and pass them in.

\needspace{14\baselineskip}
\begin{ridercode}{Java}
// Impure: loads data inside the "calculation"
public double summarizeObservations(MissionId missionId)
{
    List<ObservationEntity> rows =
        observationDao.findByMission(missionId);

    double total = 0.0;
    for (ObservationEntity row : rows)
    {
        total += row.getSignalStrength();
    }

    return total;
}
\end{ridercode}
\captionof{codesnippet}{Impure summary: database access mixed with calculation}
\label{lst:fcis-impure-summarize}

\needspace{10\baselineskip}
\begin{ridercode}{Java}
// Pure: same inputs always yield the same total
public static double summarizeSignalStrength(
    List<ObservationSnapshot> observations
)
{
    return observations.stream()
        .mapToDouble(ObservationSnapshot::signalStrength)
        .sum();
}
\end{ridercode}
\captionof{codesnippet}{Pure summary: explicit inputs, explicit result}
\label{lst:fcis-pure-summarize}

\subsection{Immutable and Mutable Data}
\label{sec:fcis_immutable_mutable}

\textbf{Mutable data} can be changed after creation; \textbf{immutable data} cannot. When several threads or several layers share mutable objects, reasoning about who changed what becomes expensive. Pure functions work best with immutable inputs: they return new values instead of altering arguments in place.

In Java, \texttt{record} types are a practical default for values crossing the core boundary. Prefer \texttt{List.copyOf} or unmodifiable lists when exposing collections from the core, and treat any object passed \textit{into} the core as read-only unless you have a compelling reason otherwise. The shell may still use mutable JDBC entities internally; map them to immutable snapshots before calling the core.

\subsection{Side Effects}
\label{sec:fcis_side_effects}

A \textbf{side effect} is any interaction with the outside world or any observable change beyond returning a value. Common side effects include:
\begin{itemize}
    \item database reads and writes,
    \item network and message-broker calls,
    \item file system access,
    \item logging and metrics,
    \item mutation of shared or global state,
    \item reading non-deterministic sources such as the system clock or a random number generator.
\end{itemize}
Side effects are necessary in real systems. FCIS confines them to the imperative shell so the functional core stays a pipeline of honest, pure transformations.

\subsection{Applying These Ideas in Java}
\label{sec:fcis_java_guidance}

Java is not a functional language, but FCIS maps cleanly onto familiar tools:
\begin{itemize}
    \item \textbf{Core package:} \texttt{static} pure methods or small \texttt{final} classes with no infrastructure imports. No Spring, JDBC, JPA, or HTTP types.
    \item \textbf{Value types:} Java \texttt{record} for inputs and outputs (\texttt{ObservationSnapshot}, \texttt{ReconciledTelemetry}, \texttt{AnnularSector}).
    \item \textbf{Imperative shell:} initially the controller or event handler; later an application service when entry points multiply. The shell loads data, calls the core, then persists or publishes results.
    \item \textbf{Immutability:} \texttt{final} fields on shell classes; defensive copies at the core boundary.
    \item \textbf{Optional:} express absence without \texttt{null}; the core returns data, the shell decides whether to throw or map to HTTP errors.
    \item \textbf{Push rule:} when you write an \texttt{if} in the shell, ask whether the condition could be computed by a pure function instead.\autocite{googleFcis2025}
\end{itemize}
The rest of this chapter applies these rules to two Abyssal Watch use cases.

\section{Why We Use Functional Core, Imperative Shell}
\label{sec:fcis_why}

\subsection{Testability Without Mocks}
\label{sec:fcis_why_testability}

The main payoff is \textbf{fast, focused unit tests}. When business rules live in pure functions, tests pass plain data structures in and assert on returned values. You do not need fake databases, stubbed HTTP clients, or mocked repositories to verify scoring logic or geometric fusion. Integration tests still matter for the shell---wiring, SQL, and transaction boundaries---but the expensive, branching logic is already proven in isolation.

\subsection{Clear Separation of Concerns}
\label{sec:fcis_why_separation}

FCIS forces a repeatable shape: \textbf{read} data, \textbf{compute} a result, \textbf{write} or publish. That read-compute-write rhythm matches how we describe use cases to non-developers (``load the readings, figure out where the Kaiju is, store the estimate'') and matches how we draw sequence diagrams. When calculation is buried inside DAO callbacks or controller actions, the story is harder to follow and harder to change.

\subsection{Safer Concurrency}
\label{sec:fcis_why_concurrency}

Pure functions do not share mutable state. If the core only transforms immutable values, you can call it from multiple threads without locks, and you can parallelise independent calculations when needed. The shell still coordinates I/O and transactions; the core stays embarrassingly parallel friendly.

\subsection{Relationship to Other Architectural Styles}
\label{sec:fcis_why_other_styles}

FCIS complements other styles taught in this book:
\begin{itemize}
    \item In \textbf{two-layered architecture} (Chapter~\ref{ch:2_layered}), business logic often lives in controllers. FCIS extracts that logic into a testable core while the controller remains the imperative shell---at least at first.
    \item In \textbf{three-layered architecture} (Chapter~\ref{ch:three_layered}), thin controllers delegate to services such as \texttt{ConfirmKaijuDetectionService}. That service layer is the \textbf{evolved} FCIS shell we introduce in Section~\ref{sec:fcis_evolution_service}.
    \item In \textbf{hexagonal architecture} (Chapter~\ref{ch:hexagonal}), delivery and infrastructure adapters define \textit{where} technology plugs in; FCIS defines \textit{what} must stay pure inside the orchestration path.
    \item FCIS does not require a service class or port interfaces on day one. A minimal adoption is one extracted pure class and a slimmer controller method.
\end{itemize}

\section{The Imperative Sandwich}
\label{sec:fcis_imperative_sandwich}

The \textbf{imperative sandwich} (also called the \textbf{functional core sandwich}) is the request-level flow that FCIS encourages:\autocite{seemannFcis,googleFcis2025}
\begin{enumerate}
    \item \textbf{Read (imperative):} the shell fetches data from databases, APIs, sensors, or user input.
    \item \textbf{Compute (functional):} the shell passes plain data into the core; the core returns a new value with no side effects.
    \item \textbf{Write (imperative):} the shell persists results, sends responses, publishes events, or triggers hardware.
\end{enumerate}
The imperative layers are the bread; the pure core is the filling. Keep the bread thin and the filling rich.

\needspace{12\baselineskip}
\tododiagram{imperative sandwich: read I/O, pure functional core, write I/O}
\begin{figure}[htbp]
\centering
\diagramplaceholder{read (shell) $\rightarrow$ compute (core) $\rightarrow$ write (shell)}
\caption{The imperative sandwich: side effects on the outside, pure logic in the middle (diagram to be added).}
\label{fig:fcis_imperative_sandwich}
\end{figure}

Figure~\ref{fig:fcis_imperative_sandwich} summarises the pattern. In the simple two-layer form described next, the ``shell'' is the delivery layer---a controller or event handler. Compare that to a \textbf{fat controller} where calculation is interleaved with SQL and mapping: the same use case works, but tests must mimic the world to reach the logic.

\needspace{10\baselineskip}
\tododiagram{fat service vs imperative sandwich side by side}
\begin{figure}[htbp]
\centering
\diagramplaceholder{left: interleaved I/O and logic in one service; right: thin shell around pure core}
\caption{Fat service versus imperative sandwich structure (diagram to be added).}
\label{fig:fcis_fat_vs_sandwich}
\end{figure}

\subsection{Sandwich in Pseudocode}
\label{sec:fcis_sandwich_pseudocode}

The following pseudocode (adapted from common FCIS presentations) shows the rhythm for processing an order status. The Java examples later follow the same three steps.

\needspace{14\baselineskip}
\begin{ridercode}{Java}
// Imperative shell
void processOrder(OrderId orderId)
{
    Order order = db.find(orderId);

    OrderStatus newStatus = calculateNextStatus(
        order.status(),
        order.amount()
    );

    db.updateStatus(orderId, newStatus);
}

// Functional core
OrderStatus calculateNextStatus(
    OrderStatus status,
    double amount
)
{
    if (status == PENDING && amount > 100.0)
    {
        return REVIEW_REQUIRED;
    }

    return APPROVED;
}
\end{ridercode}
\captionof{codesnippet}{Imperative sandwich in pseudocode}
\label{lst:fcis-sandwich-pseudocode}

Do not pass heavy domain objects with hidden behaviour into the core if you can avoid it. Prefer simple records or maps that the core can treat as dumb data.\autocite{fritzscheFcis} The shell is responsible for mapping from persistence entities and framework types into those records.

\section{Starting Simple: Two Layers}
\label{sec:fcis_two_layers}

FCIS does not require a service class on day one. When a use case has \textbf{one entry point}---a single REST endpoint, one JavaFX screen, or one CLI command---the \textbf{controller or handler is the imperative shell}. It performs the sandwich: load data, call the pure core, then persist or respond.

\needspace{12\baselineskip}
\tododiagram{two-layer FCIS: Web API or JavaFX delivery layer → functional core}
\begin{figure}[htbp]
\centering
\diagramplaceholder{delivery (impure) $\rightarrow$ core (pure)}
\caption{Minimal FCIS: one impure delivery layer and one pure core (diagram to be added).}
\label{fig:fcis_two_layers}
\end{figure}

Figure~\ref{fig:fcis_two_layers} shows the starting point for this chapter. The delivery class may inject DAOs, sensor drivers, or file readers directly. That couples infrastructure to the entry point, which is acceptable for a small system with one way to trigger each use case. What matters is that \textbf{calculation} lives in the core, not that every folder name is perfect.

This is an improvement over the fat-controller risk from Chapter~\ref{ch:2_layered}: the controller may still own I/O, but it delegates rules and transformations to honest pure functions. Sensor polling, JDBC, and HTTP response mapping are shell work regardless of whether they sit in a controller or a service.

When the same behaviour must be reachable from a Web API, a CLI job, and a GUI button, we promote the sandwich into a shared \textbf{application service} and thin out the delivery adapters. Section~\ref{sec:fcis_evolution_service} shows that step using SCN-04.

\section{Example: Triangulate Kaiju Position (SCN-03)}
\label{sec:fcis_example_triangulation}

This example revisits the triangulation use case from Chapter~\ref{ch:hexagonal}---our \textbf{known} example. We suspect a Kaiju has left a Breach; multiple sensors in the area return uncertain readings. Each reading is an \textit{annular sector}---a direction and distance range, not a single point. We fuse several sectors to narrow the likely position to a small \texttt{EstimatedArea} polygon.

In the hexagonal chapter, ports and adapters framed the problem; fusion math remained a private stub inside \texttt{TriangulateKaijuPositionService}. Here we start \textbf{one level simpler}: extract that stub into \texttt{KaijuPositionEstimator} and let a \texttt{TriangulationController} act as the imperative shell. Section~\ref{sec:fcis_evolution_service} explains when to promote the same pattern into a dedicated service.

\subsection{Use Case Summary}
\label{sec:fcis_triangulation_use_case}

The input is a suspected \texttt{Location} and a search \texttt{radius}. The output is an \texttt{EstimatedArea}: a list of \texttt{Location} points outlining a polygon. Sensors within the radius are loaded, polled through type-specific drivers, and converted to \texttt{AnnularSector} values. The geometric fusion is pure mathematics on those sectors; everything involving the ocean, hardware, or database is shell work.

%\needspace{20\baselineskip}
%\begin{figure}[htbp]
%    \centering
%    \includegraphics[width=0.5\textwidth]{assets/annular-sector-triangulation.png}
%    \caption{Triangulation of multiple annular sectors (reused from the hexagonal chapter).}
%    \label{fig:fcis_annular_triangulation}
%\end{figure}

\begin{figure}[htbp]
    \centering
    % Placeholder for triangulation diagram
    %\includegraphics[width=0.6\textwidth]{assets/triangulation-placeholder.png}
    \fbox{\parbox{0.6\textwidth}{\centering Placeholder: triangulation of annular sectors\\[2ex]\strut}}
    \caption{Triangulation of multiple annular sectors to estimate Kaiju position.}
    \label{fig:fcis_annular_triangulation}
\end{figure}
\todo{Replace placeholder with a domain-specific triangulation diagram illustrating the intersection of annular sectors. Reference annular sector diagram from hexagonal chapter.}


Figure~\ref{fig:fcis_annular_triangulation} illustrates the idea: each sensor contributes an uncertain region; intersection narrows the Kaiju's likely position. For annular-sector terminology, see Chapter~\ref{ch:hexagonal} and Figure~\ref{fig:annular-sector} there.

\subsection{The Problem: Logic Inside the Service}
\label{sec:fcis_triangulation_problem}

Recall Code snippet~\ref{lst:triangulate-service-compute-area} from Chapter~\ref{ch:hexagonal}: \texttt{computeEstimatedArea} is pure logic trapped inside an impure service class. It is hard to unit test without constructing the whole service and mocking three ports, even though the geometry only needs a \texttt{List<AnnularSector>}.

The FCIS refactor extracts that helper into a dedicated \textbf{functional core} class and leaves a \textbf{thin imperative shell} at the delivery boundary.

\subsection{Functional Core: KaijuPositionEstimator}
\label{sec:fcis_triangulation_core}

We place pure types and logic under \texttt{com.abyssalwatch.core.triangulation}. \texttt{AnnularSector} is a record describing one sensor's uncertain region relative to a target. \texttt{KaijuPositionEstimator} exposes a single honest entry point.

\needspace{16\baselineskip}
\begin{ridercode}{Java}
public record AnnularSector(
    Location sensorLocation,
    double bearingDegrees,
    double minDistanceMetres,
    double maxDistanceMetres,
    double spreadDegrees
)
{
}

public final class KaijuPositionEstimator
{
    private KaijuPositionEstimator()
    {
    }

    public static List<Location> estimateArea(
        List<AnnularSector> sectors
    )
    {
        if (sectors.isEmpty())
        {
            return List.of();
        }

        // Simplified stub: real fusion would intersect sectors.
        return sectors.stream()
            .map(KaijuPositionEstimator::centrePoint)
            .distinct()
            .toList();
    }

    private static Location centrePoint(AnnularSector sector)
    {
        double midDistance =
            (sector.minDistanceMetres()
                + sector.maxDistanceMetres()) / 2.0;

        return Geo.offset(
            sector.sensorLocation(),
            sector.bearingDegrees(),
            midDistance
        );
    }
}
\end{ridercode}
\captionof{codesnippet}{\texttt{KaijuPositionEstimator} functional core (simplified fusion)}
\label{lst:fcis-kaiju-position-estimator}

The implementation is deliberately a stub---the book is about architecture, not computational geometry---but the \textit{boundary} is real: no database, no sensor driver, no logging. Production code would replace the body of \texttt{estimateArea} with a proper intersection algorithm while keeping the signature unchanged.

\subsection{Imperative Shell: TriangulationController}
\label{sec:fcis_triangulation_shell}

For a single Web API entry point, the controller is the imperative shell. It lives under \texttt{com.abyssalwatch.presentation.web} and injects infrastructure collaborators directly---a deliberate simplification for stage one. A JavaFX button handler could follow the same sandwich shape.

\needspace{28\baselineskip}
\begin{ridercode}{Java}
@RestController
@RequestMapping("/triangulation")
public final class TriangulationController
{
    private final SensorRepository sensorRepository;
    private final SensorDriverFactory sensorDriverFactory;

    public TriangulationController(
        SensorRepository sensorRepository,
        SensorDriverFactory sensorDriverFactory
    )
    {
        this.sensorRepository = sensorRepository;
        this.sensorDriverFactory = sensorDriverFactory;
    }

    @PostMapping
    public EstimatedArea triangulate(
        @RequestBody TriangulationRequest request
    )
    {
        List<Sensor> sensors =
            sensorRepository.findAllByCircle(
                request.location(),
                request.radius()
            );

        List<AnnularSector> sectors =
            pollSensors(sensors, request.location());

        List<Location> polygon =
            KaijuPositionEstimator.estimateArea(sectors);

        return new EstimatedArea(polygon);
    }

    private List<AnnularSector> pollSensors(
        List<Sensor> sensors,
        Location target
    )
    {
        List<AnnularSector> sectors = new ArrayList<>();

        for (Sensor sensor : sensors)
        {
            SensorDriver driver =
                sensorDriverFactory.create(sensor.type());

            sectors.add(driver.poll(target));
        }

        return List.copyOf(sectors);
    }
}
\end{ridercode}
\captionof{codesnippet}{\texttt{TriangulationController} as imperative shell}
\label{lst:fcis-triangulate-shell}

Read the endpoint top to bottom as an imperative sandwich: \texttt{findAllByCircle} and \texttt{pollSensors} are reads; \texttt{KaijuPositionEstimator.estimateArea} is the pure compute step; returning \texttt{EstimatedArea} is the response to the caller (persisting the estimate, if needed, would be another shell concern).

\subsection{Testing the Core}
\label{sec:fcis_triangulation_test}

Tests target the core directly. No \texttt{SensorRepository} fake is required to prove that three sectors produce a non-empty polygon.

\needspace{18\baselineskip}
\begin{ridercode}{Java}
class KaijuPositionEstimatorTest
{
    @Test
    void estimateArea_withThreeSectors_returnsNonEmptyPolygon()
    {
        Location origin = new Location(
            new Latitude(0.0),
            new Longitude(0.0)
        );

        List<AnnularSector> sectors = List.of(
            sector(origin, 45.0, 100.0, 200.0),
            sector(origin, 90.0, 100.0, 200.0),
            sector(origin, 135.0, 100.0, 200.0)
        );

        List<Location> polygon =
            KaijuPositionEstimator.estimateArea(sectors);

        assertFalse(polygon.isEmpty());
    }

    private static AnnularSector sector(
        Location sensor,
        double bearing,
        double minDist,
        double maxDist
    )
    {
        return new AnnularSector(
            sensor,
            bearing,
            minDist,
            maxDist,
            5.0
        );
    }
}
\end{ridercode}
\captionof{codesnippet}{Unit test for \texttt{KaijuPositionEstimator} without mocks}
\label{lst:fcis-triangulation-test}

If fusion rules change---new weighting for thermal sensors, for example---you edit \texttt{KaijuPositionEstimator} and its tests. Sensor driver protocols can change independently in the controller. When a CLI batch job also needs coordinate lock, extract the sandwich into a service as in Section~\ref{sec:fcis_evolution_service}.

\section{Example: Reconcile Mission Telemetry (SCN-04)}
\label{sec:fcis_example_reconciliation}

This is our \textbf{new} example. Scout Vessel \texttt{S-402 ``Mantis''} had intermittent uplink during the \texttt{Retrieval} phase of a mission. Some observations synced over the air; the rest arrive when the vessel docks at the Command Carrier. Before the mission can move toward \texttt{Completed} (invariant INV-05), the system must produce one consistent observation timeline.

That merge-and-deduplicate pipeline is ideal FCIS material: load both batches in the shell, reconcile in the core, persist the result in the shell.

\subsection{Use Case Summary}
\label{sec:fcis_reconciliation_use_case}

When a retrieval order closes the tracking window, operators need a single authoritative observation set. The workflow is:
\begin{enumerate}
    \item \textbf{Read:} load observations already synced for the mission, plus the dock-upload batch from the returning Scout Vessel,
    \item \textbf{Compute:} merge, sort by timestamp, and deduplicate overlapping readings,
    \item \textbf{Write:} persist the reconciled timeline so downstream mission-closure checks can proceed.
\end{enumerate}

\needspace{12\baselineskip}
\tododiagram{telemetry reconciliation: stored batch + dock upload → TelemetryReconciler → persist}
\begin{figure}[htbp]
\centering
\diagramplaceholder{stored batch + dock upload $\rightarrow$ TelemetryReconciler $\rightarrow$ persist}
\caption{Telemetry reconciliation data flow (diagram to be added).}
\label{fig:fcis_reconciliation_flow}
\end{figure}

Figure~\ref{fig:fcis_reconciliation_flow} shows the imperative sandwich: two read paths feed one pure merge; the shell owns persistence and any mission-state transition.

\subsection{Functional Core: TelemetryReconciler}
\label{sec:fcis_reconciliation_core}

Core types live under \texttt{com.abyssalwatch.core.telemetry}. \texttt{ObservationSnapshot} is the immutable input shape; the shell maps JDBC entities and dock-upload records to snapshots before calling the core.

\needspace{28\baselineskip}
\begin{ridercode}{Java}
public record ObservationSnapshot(
    ObservationId id,
    Instant observedAt,
    SensorType sensorType,
    double signalStrength
)
{
}

public record ReconciliationPolicy(
    Duration duplicateWindow
)
{
}

public record ReconciledTelemetry(
    List<ObservationSnapshot> observations,
    int duplicateCount,
    double completenessRatio
)
{
}

public final class TelemetryReconciler
{
    private TelemetryReconciler()
    {
    }

    public static ReconciledTelemetry merge(
        List<ObservationSnapshot> alreadySynced,
        List<ObservationSnapshot> dockUpload,
        ReconciliationPolicy policy
    )
    {
        List<ObservationSnapshot> combined =
            new ArrayList<>(alreadySynced);
        combined.addAll(dockUpload);

        combined.sort(Comparator.comparing(
            ObservationSnapshot::observedAt
        ));

        List<ObservationSnapshot> deduped =
            deduplicate(combined, policy.duplicateWindow());

        int duplicateCount =
            combined.size() - deduped.size();

        int expectedCount = combined.size();
        double completeness = expectedCount == 0
            ? 1.0
            : (double) deduped.size() / expectedCount;

        return new ReconciledTelemetry(
            List.copyOf(deduped),
            duplicateCount,
            completeness
        );
    }

    private static List<ObservationSnapshot> deduplicate(
        List<ObservationSnapshot> sorted,
        Duration window
    )
    {
        List<ObservationSnapshot> result =
            new ArrayList<>();

        for (ObservationSnapshot current : sorted)
        {
            Optional<ObservationSnapshot> existing =
                findDuplicate(result, current, window);

            if (existing.isEmpty())
            {
                result.add(current);
                continue;
            }

            if (current.signalStrength()
                > existing.get().signalStrength())
            {
                result.remove(existing.get());
                result.add(current);
            }
        }

        return result;
    }

    private static Optional<ObservationSnapshot> findDuplicate(
        List<ObservationSnapshot> kept,
        ObservationSnapshot candidate,
        Duration window
    )
    {
        for (ObservationSnapshot other : kept)
        {
            boolean sameSensor = other.sensorType()
                .equals(candidate.sensorType());

            Duration gap = Duration.between(
                other.observedAt(),
                candidate.observedAt()
            ).abs();

            if (sameSensor && gap.compareTo(window) <= 0)
            {
                return Optional.of(other);
            }
        }

        return Optional.empty();
    }
}
\end{ridercode}
\captionof{codesnippet}{\texttt{TelemetryReconciler} functional core}
\label{lst:fcis-reconciliation-core}

The deduplication rule keeps the stronger signal when two readings share a sensor type within the configured window. Production code might key on \texttt{ObservationId} or sensor serial numbers; here the focus is the architectural boundary: no DAO, no file reader, no mission-state update inside the core.

\subsection{Imperative Shell: MissionTelemetryController}
\label{sec:fcis_reconciliation_shell}

With one REST endpoint for dock-side reconciliation, the controller is again the imperative shell. It loads both batches, maps to snapshots, calls the reconciler, and persists the merged result.

\needspace{30\baselineskip}
\begin{ridercode}{Java}
@RestController
@RequestMapping("/missions")
public final class MissionTelemetryController
{
    private final ObservationDao observationDao;
    private final DockUploadReader dockUploadReader;
    private final UnitOfWork unitOfWork;

    @PostMapping("/{missionId}/reconcile")
    public ReconciledTelemetry reconcile(
        @PathVariable MissionId missionId,
        @RequestBody ReconcileTelemetryRequest request
    )
    {
        List<ObservationSnapshot> synced =
            loadSynced(missionId);

        List<ObservationSnapshot> dockBatch =
            dockUploadReader.read(request.vesselId());

        ReconciledTelemetry merged =
            TelemetryReconciler.merge(
                synced,
                dockBatch,
                request.policy()
            );

        persistMerged(missionId, merged);

        return merged;
    }

    private List<ObservationSnapshot> loadSynced(
        MissionId missionId
    )
    {
        return observationDao
            .findSyncedByMission(missionId)
            .stream()
            .map(this::toSnapshot)
            .toList();
    }

    private ObservationSnapshot toSnapshot(
        ObservationEntity row
    )
    {
        return new ObservationSnapshot(
            new ObservationId(row.getId()),
            row.getObservedAt(),
            new SensorType(row.getSensorType()),
            row.getSignalStrength()
        );
    }

    private void persistMerged(
        MissionId missionId,
        ReconciledTelemetry merged
    )
    {
        unitOfWork.execute(() ->
        {
            observationDao.replaceMissionObservations(
                missionId,
                merged.observations()
            );
            return null;
        });
    }
}
\end{ridercode}
\captionof{codesnippet}{\texttt{MissionTelemetryController} as imperative shell}
\label{lst:fcis-reconciliation-shell}

Whether the mission may transition toward \texttt{Completed} is a separate shell concern for now: the controller could read \texttt{ReconciledTelemetry.completenessRatio()} and vessel recovery state, then apply INV-05. Section~\ref{sec:fcis_evolution_service} moves this orchestration into a shared service when more entry points appear.

\subsection{Testing the Core with Table-Driven Cases}
\label{sec:fcis_reconciliation_test}

Table-driven tests document reconciliation behaviour without a database.

\needspace{24\baselineskip}
\begin{ridercode}{Java}
class TelemetryReconcilerTest
{
    private static final ReconciliationPolicy POLICY =
        new ReconciliationPolicy(Duration.ofMinutes(5));

    @Test
    void merge_duplicateWithinWindow_keepsStrongerSignal()
    {
        Instant t = Instant.parse("2026-06-01T12:00:00Z");

        List<ObservationSnapshot> synced = List.of(
            snap("obs-1", t, "Radiation", 30.0)
        );

        List<ObservationSnapshot> dock = List.of(
            snap("obs-2", t.plusSeconds(60), "Radiation", 45.0)
        );

        ReconciledTelemetry result =
            TelemetryReconciler.merge(synced, dock, POLICY);

        assertEquals(1, result.observations().size());
        assertEquals(45.0, result.observations().get(0)
            .signalStrength());
        assertEquals(1, result.duplicateCount());
    }

    @Test
    void merge_dockOnlyRecords_appended()
    {
        Instant t = Instant.parse("2026-06-01T12:00:00Z");

        List<ObservationSnapshot> synced = List.of(
            snap("obs-1", t, "Radiation", 30.0)
        );

        List<ObservationSnapshot> dock = List.of(
            snap("obs-2", t.plusHours(1), "Acoustic", 20.0)
        );

        ReconciledTelemetry result =
            TelemetryReconciler.merge(synced, dock, POLICY);

        assertEquals(2, result.observations().size());
    }

    @Test
    void merge_emptyDock_returnsSyncedUnchanged()
    {
        Instant t = Instant.parse("2026-06-01T12:00:00Z");

        List<ObservationSnapshot> synced = List.of(
            snap("obs-1", t, "Radiation", 30.0)
        );

        ReconciledTelemetry result =
            TelemetryReconciler.merge(
                synced,
                List.of(),
                POLICY
            );

        assertEquals(1, result.observations().size());
        assertEquals(0, result.duplicateCount());
    }

    private static ObservationSnapshot snap(
        String id,
        Instant time,
        String sensorType,
        double strength
    )
    {
        return new ObservationSnapshot(
            new ObservationId(id),
            time,
            new SensorType(sensorType),
            strength
        );
    }
}
\end{ridercode}
\captionof{codesnippet}{Table-driven tests for telemetry reconciliation}
\label{lst:fcis-reconciliation-test}

Changing the deduplication window or the keep-stronger rule is a one-file change with tests that run in milliseconds.\todoinref{property-based testing chapter for fuzzing merged ObservationSnapshot lists}

\section{When Entry Points Multiply: Adding a Service Layer}
\label{sec:fcis_evolution_service}

The two-layer shape is a good start, but it does not scale when the \textbf{same use case} must be triggered from several places. Suppose SCN-04 reconciliation must run from:
\begin{itemize}
    \item the REST endpoint in Code snippet~\ref{lst:fcis-reconciliation-shell},
    \item a CLI import when a dock upload file lands on the Command Carrier,
    \item and a nightly batch job that catches missions stuck in \texttt{Retrieval}.
\end{itemize}
Copying the full imperative sandwich into three classes duplicates I/O sequencing and drifts out of sync. The fix is to extract a shared \textbf{application service} and leave delivery classes thin.

\subsection{The Problem: Duplicated Sandwiches}
\label{sec:fcis_evolution_problem}

Each new entry point would repeat: load synced observations, read the dock batch, call \texttt{TelemetryReconciler.merge}, persist. The functional core does not change---only the duplication in impure orchestration grows. This is the same pressure that led to application services in Chapter~\ref{ch:three_layered}, where \texttt{KaijuController} delegates to \texttt{ConfirmKaijuDetectionService}.

\subsection{Extract the Application Service}
\label{sec:fcis_evolution_extract}

Move the sandwich from \texttt{MissionTelemetryController} into \texttt{ReconcileMissionTelemetryService}. The service owns read--compute--write; \texttt{TelemetryReconciler} is unchanged.

\needspace{26\baselineskip}
\begin{ridercode}{Java}
public final class ReconcileMissionTelemetryService
{
    private final ObservationDao observationDao;
    private final DockUploadReader dockUploadReader;
    private final UnitOfWork unitOfWork;

    public ReconciledTelemetry reconcile(
        MissionId missionId,
        VesselId vesselId,
        ReconciliationPolicy policy
    )
    {
        List<ObservationSnapshot> synced =
            loadSynced(missionId);

        List<ObservationSnapshot> dockBatch =
            dockUploadReader.read(vesselId);

        ReconciledTelemetry merged =
            TelemetryReconciler.merge(
                synced,
                dockBatch,
                policy
            );

        persistMerged(missionId, merged);

        return merged;
    }

    // loadSynced, toSnapshot, persistMerged unchanged
    // from Code snippet~\ref{lst:fcis-reconciliation-shell}
}
\end{ridercode}
\captionof{codesnippet}{\texttt{ReconcileMissionTelemetryService} shared orchestration shell}
\label{lst:fcis-reconciliation-service}

The private helpers from Code snippet~\ref{lst:fcis-reconciliation-shell} move here unchanged. The functional core did not need a single edit.

\subsection{Split Delivery from Infrastructure}
\label{sec:fcis_evolution_split}

With a service in place, impure code falls into two roles:
\begin{itemize}
    \item \textbf{Delivery adapters} translate an entry point into a service call. They map HTTP bodies, CLI arguments, or JavaFX form values, and map service results back to responses. They should not query the database directly.
    \item \textbf{Infrastructure} provides outbound capabilities the service needs: \texttt{ObservationDao}, \texttt{DockUploadReader}, message brokers, sensor drivers. The service depends on infrastructure; delivery depends on the service.
\end{itemize}

These roles mirror ideas from layered and hexagonal chapters, but the names here are deliberately plain. FCIS still cares only about \textbf{where purity lives}: inside \texttt{TelemetryReconciler}, not in delivery or infrastructure.

\needspace{12\baselineskip}
\tododiagram{delivery → application service → core; service → infrastructure}
\begin{figure}[htbp]
\centering
\diagramplaceholder{delivery $\rightarrow$ service $\rightarrow$ core; service $\rightarrow$ infrastructure}
\caption{Evolved FCIS: thin delivery, service as orchestration shell, infrastructure for I/O (diagram to be added).}
\label{fig:fcis_evolution_layers}
\end{figure}

The REST controller becomes thin:

\needspace{14\baselineskip}
\begin{ridercode}{Java}
@RestController
@RequestMapping("/missions")
public final class MissionTelemetryController
{
    private final ReconcileMissionTelemetryService service;

    @PostMapping("/{missionId}/reconcile")
    public ReconciledTelemetry reconcile(
        @PathVariable MissionId missionId,
        @RequestBody ReconcileTelemetryRequest request
    )
    {
        return service.reconcile(
            missionId,
            request.vesselId(),
            request.policy()
        );
    }
}
\end{ridercode}
\captionof{codesnippet}{Thin delivery adapter after service extraction}
\label{lst:fcis-reconciliation-controller-thin}

A CLI handler follows the same pattern: parse arguments, call \texttt{service.reconcile(...)}, print or log the result. Triangulation (SCN-03) promotes the same way when a batch job or second API also needs coordinate lock.

\subsection{What Stayed Pure}
\label{sec:fcis_evolution_pure_unchanged}

\texttt{TelemetryReconciler} and its unit tests are identical before and after the refactor. That stability is the payoff: delivery and infrastructure arrangements can evolve while the rules under test stay fixed. Promote to a service when entry points multiply; do not promote earlier just because other chapters use service classes.

\section{Recommended Package Layout}
\label{sec:fcis_package_layout}

FCIS benefits from explicit package boundaries. The invariant is always: \texttt{core} must not depend on delivery, infrastructure, or persistence types.

\subsection{Simple Layout: Core and Delivery}
\label{sec:fcis_package_layout_simple}

Stage one maps cleanly to three roots: \texttt{core}, \texttt{presentation}, and \texttt{persistence}. Controllers in \texttt{presentation} call the core and infrastructure directly.

\needspace{16\baselineskip}
\begin{figure}[htbp]
\centering
\begin{minipage}{0.85\textwidth}
\dirtree{%
.1 \foldericon\ com.abyssalwatch/.
.2 \foldericon\ core/.
.3 \foldericon\ triangulation/.
.4 \fileicon\ KaijuPositionEstimator.java.
.4 \fileicon\ AnnularSector.java.
.3 \foldericon\ telemetry/.
.4 \fileicon\ TelemetryReconciler.java.
.4 \fileicon\ ReconciledTelemetry.java.
.4 \fileicon\ ObservationSnapshot.java.
.2 \foldericon\ presentation/.
.3 \foldericon\ web/.
.4 \fileicon\ TriangulationController.java.
.4 \fileicon\ MissionTelemetryController.java.
.2 \foldericon\ persistence/.
.3 \fileicon\ JdbcObservationDao.java.
.3 \fileicon\ JdbcSensorDao.java.
}
\end{minipage}
\caption{Simple FCIS layout: core plus delivery controllers with direct infrastructure access.}
\label{fig:fcis_package_layout_simple}
\end{figure}

In this layout, map \texttt{ObservationEntity} to \texttt{ObservationSnapshot} in the controller before calling \texttt{TelemetryReconciler}. \texttt{Sensor} records and driver polling stay in the controller for SCN-03.

\subsection{Evolved Layout: Delivery, Service, and Infrastructure}
\label{sec:fcis_package_layout_evolved}

After Section~\ref{sec:fcis_evolution_service}, add \texttt{application} for orchestration services and move persistence and file readers under \texttt{infrastructure}. Delivery adapters stay thin.

\needspace{16\baselineskip}
\begin{figure}[htbp]
\centering
\begin{minipage}{0.85\textwidth}
\dirtree{%
.1 \foldericon\ com.abyssalwatch/.
.2 \foldericon\ core/.
.3 \foldericon\ telemetry/.
.4 \fileicon\ TelemetryReconciler.java.
.2 \foldericon\ application/.
.3 \foldericon\ telemetry/.
.4 \fileicon\ ReconcileMissionTelemetryService.java.
.2 \foldericon\ delivery/.
.3 \foldericon\ web/.
.4 \fileicon\ MissionTelemetryController.java.
.3 \foldericon\ cli/.
.4 \fileicon\ ReconcileTelemetryCommand.java.
.2 \foldericon\ infrastructure/.
.3 \foldericon\ persistence/.
.4 \fileicon\ JdbcObservationDao.java.
.3 \foldericon\ files/.
.4 \fileicon\ DockUploadReader.java.
}
\end{minipage}
\caption{Evolved FCIS layout: delivery, application service, and infrastructure.}
\label{fig:fcis_package_layout_evolved}
\end{figure}

Here the entity-to-snapshot mapping moves to the application service boundary. Delivery only maps HTTP or CLI types to service parameters.

\section{Testing the Functional Core}
\label{sec:fcis_testing}

\subsection{What to Test Where}
\label{sec:fcis_testing_where}

\begin{itemize}
    \item \textbf{Functional core:} unit tests with fixed inputs; no Spring context, no test containers, no mocked repositories. Aim for branch coverage on rules and calculations.
    \item \textbf{Simple imperative shell:} integration tests through the controller verify wiring---SQL mappings, transaction rollback, HTTP status codes.
    \item \textbf{Evolved imperative shell:} test the application service with fake infrastructure collaborators; keep delivery adapter tests minimal.
    \item \textbf{End-to-end:} optional smoke tests through the UI or API; not the primary place to prove reconciliation logic.
\end{itemize}

\subsection{Contrast with Port-Based Testing}
\label{sec:fcis_testing_vs_ports}

Hexagonal architecture (Chapter~\ref{ch:hexagonal}) tests use cases by substituting fake outbound ports. That remains valuable for service-level orchestration after Section~\ref{sec:fcis_evolution_service}. FCIS reduces how much logic those fakes must exercise: if \texttt{KaijuPositionEstimator} is already proven, a controller or service integration test only needs to assert that the shell calls the estimator and maps the result---not that intersection math is correct.

\subsection{Guidelines}
\label{sec:fcis_testing_guidelines}

\begin{itemize}
    \item Pass time, configuration thresholds, and feature flags \textit{into} the core as parameters rather than reading static singletons.
    \item Prefer many small pure functions over one large \texttt{merge} method when branches multiply; each function gets its own focused test.
    \item When a shell test fails, determine whether the bug is in wiring (shell) or rules (core) before fixing; mixing fixes across the boundary recreates a fat service.
    \item Reconciliation cores such as \texttt{TelemetryReconciler} are good candidates for property-based tests later.\todoinref{property-based testing chapter for merge invariants}
\end{itemize}

\section{Practical Trade-offs}
\label{sec:fcis_tradeoffs}

Functional Core, Imperative Shell is a deliberate way to spend a little structure up front in exchange for testable calculations. It is not free, and it is not mandatory for every class.

\subsection{Benefits}
\label{sec:fcis_pros}

When the investment fits the problem, FCIS delivers:
\begin{itemize}
    \item \textbf{Fast rule tests:} complex logic runs in unit tests without mocks, as in Code snippets~\ref{lst:fcis-triangulation-test} and~\ref{lst:fcis-reconciliation-test},
    \item \textbf{Maintainability:} business rule changes concentrate in \texttt{core} packages, away from JDBC and HTTP,
    \item \textbf{Readability:} the imperative sandwich makes read-compute-write explicit, as in Figure~\ref{fig:fcis_imperative_sandwich},
    \item \textbf{Concurrency safety:} pure cores avoid shared mutable state during calculation,
    \item \textbf{Composable logic:} pure functions can be reused across use cases, batch jobs, and reporting pipelines without copying I/O code,
    \item and \textbf{Alignment with other styles:} FCIS sharpens hexagonal and layered boundaries rather than replacing them, as discussed in Section~\ref{sec:fcis_why_other_styles}.
\end{itemize}

\subsection{Costs}
\label{sec:fcis_cons}

Those benefits come with costs teams should expect:
\begin{itemize}
    \item \textbf{Extra types and mapping:} snapshots and DTOs at the core boundary add classes and conversion code,
    \item \textbf{Split methods:} one procedural controller or service becomes a shell class plus one or more core classes,
    \item \textbf{Timing the service promotion:} extracting a service too early adds folders before you need them; delaying too long duplicates sandwiches across delivery adapters,
    \item \textbf{Discipline required:} without import rules, \texttt{core} packages accumulate static database access and ``temporary'' logging,
    \item \textbf{Not every line benefits:} pass-through CRUD with no rules gains little from FCIS,
    \item and \textbf{Learning curve:} developers must recognise what belongs in the core versus the shell.
\end{itemize}

\subsection{Things to Consider Before Deciding}
\label{sec:fcis_when_to_choose}

Favour Functional Core, Imperative Shell when:
\begin{itemize}
    \item calculations, scoring, reconciliation, or multi-step data transformations dominate a feature,
    \item rules change frequently and must be regression-tested,
    \item the same logic must run in the live app, batch jobs, and offline analysis,
    \item you need deterministic tests without standing up infrastructure,
    \item concurrent or parallel execution of business logic is on the roadmap,
    \item and existing services are hard to test because logic is interleaved with SQL and messaging.
\end{itemize}

Question or defer FCIS when:
\begin{itemize}
    \item the feature is thin CRUD with at most one or two trivial branches,
    \item the team will not enforce package or import boundaries,
    \item you are building a throwaway prototype,
    \item all behaviour is already a single honest pure function with no I/O in the same class,
    \item and splitting would produce more boilerplate than testable logic.
\end{itemize}

When uncertain, extract one pure function from the hottest path---often a private helper in an oversized service---and measure whether tests become simpler. Deepen the split only where the payoff is clear.

\subsection{Common Implementation Mistakes}
\label{sec:fcis_common_mistakes}

The first group of mistakes weakens purity in the core:
\begin{itemize}
    \item logging or metrics inside core methods,
    \item injecting repositories or HTTP clients into core classes,
    \item reading configuration singletons or environment variables in the core instead of passing values as parameters,
    \item mutating input lists or records passed from the shell,
    \item and calling \texttt{Instant.now()} inside the core rather than accepting time as an argument.
\end{itemize}

The next group leaves too much logic in the shell or duplicates it:
\begin{itemize}
    \item long \texttt{if/else} chains in controllers or services that should be pure functions,
    \item copying the full imperative sandwich into every delivery adapter instead of extracting an application service,
    \item reconciliation logic duplicated across REST controllers and batch importers instead of shared core code,
    \item returning different DTO shapes from the core (the core should return domain-neutral values; the shell maps to API contracts),
    \item and ``anemic shell'' services that only delegate while the core performs I/O through static helpers---a fake split.
\end{itemize}

Finally, some mistakes come from structure without discipline:
\begin{itemize}
    \item promoting to an application service before a second entry point exists,
    \item creating \texttt{core} and delivery packages but allowing cross-imports in both directions,
    \item passing live domain entities with behaviour into the core when a snapshot record would suffice,
    \item testing only through HTTP while complex math remains untested at the unit level,
    \item and applying FCIS to every class uniformly instead of focusing on calculation-heavy areas.
\end{itemize}

Most mistakes boil down to the same rule: the core must be honest and pure; the shell owns the messy world. When that holds, both Abyssal Watch examples in this chapter stay easy to evolve---deduplication rules in \texttt{TelemetryReconciler}, better fusion in \texttt{KaijuPositionEstimator}, without rewiring the ocean.

\needspace{8\baselineskip}
\subsection{Closing Thoughts}
\label{sec:fcis_closing}

We opened with functional programming vocabulary---pure and impure functions, honest signatures, immutable values, and side effects---then applied those ideas through the imperative sandwich. We started with the \textbf{two-layer} shape: controllers as imperative shells around \texttt{KaijuPositionEstimator} and \texttt{TelemetryReconciler}. Section~\ref{sec:fcis_evolution_service} showed when to add an application service and split delivery from infrastructure---without changing the functional core.

FCIS works in Java without adopting a new framework. Records, static pure methods, and disciplined imports are enough to start at the controller boundary. Add services and infrastructure folders when entry points multiply, not before. Combine FCIS with the architectural styles elsewhere in this book: layers and hexagons tell you how to organise modules; FCIS tells you how to keep the behaviour inside them testable and honest.
